\chapter{Roadmap}
\label{ch:roadmap}

\section{Basis}

This roadmap is reconstructed from three sources in the repository: the
extension README's ``Roadmap'' checklist, the extension CHANGELOG's
\texttt{[Unreleased]} section, and gaps established elsewhere in this document.
Where an item is drawn from the repository it is labelled accordingly; items
this analysis proposes are labelled RECOMMENDED. Nothing in this chapter should
be read as a commitment by the authors.

\section{Repository-stated roadmap}

The extension README lists these items as outstanding (unchecked):

\begin{longtable}{p{1.4cm}p{5.4cm}p{2.4cm}p{5.2cm}}
\caption{Items the repository itself lists as future work}\\
\toprule
\textbf{ID} & \textbf{Item} & \textbf{Source} & \textbf{Reality check} \\
\midrule

RM-01 & \texttt{expo-snips update} --- update installed screens to the latest
recipe versions. & Extension README & \textbf{Already implemented} in the CLI
(\texttt{commands/update.ts}). The README checklist is stale. \\
RM-02 & \texttt{expo-snips status} --- inspect installed screens and available
updates. & Extension README & \textbf{Already implemented} in the CLI
(\texttt{commands/status.ts}). Stale. \\
RM-03 & Tamagui styling support. & README & Genuinely outstanding; appears in the
schema and settings but has no recipe. \\
RM-04 & Vector Icons support. & README & Genuinely outstanding; enum exists, no
recipe uses it. \\
RM-05 & Custom UI primitive support. & README & Genuinely outstanding; enum
exists, no recipe uses it. \\
RM-06 & Expanded screen catalog. & README & Ongoing by nature. \\
RM-07 & Testing infrastructure. & README & Outstanding. See
Chapter~\ref{ch:testing}. \\
RM-08 & CI/CD pipeline for the VS Code extension. & README & Outstanding. See
Chapter~\ref{ch:cicd}. \\
\bottomrule
\end{longtable}

\textbf{Note:} RM-01 and RM-02 are already delivered, so the README roadmap
should be corrected. This is a small but user-visible inaccuracy.

\section{Near-term recommendations}

\textbf{Status:} RECOMMENDED. These are the items with the highest ratio of
value to effort, derived from the findings chapters.

\begin{longtable}{p{1.4cm}p{5.2cm}p{3cm}p{4.8cm}}
\caption{Near-term recommendations}\\
\toprule
\textbf{ID} & \textbf{Recommendation} & \textbf{Addresses} & \textbf{Rationale} \\
\midrule

RN-01 & Add a pull-request CI workflow that runs recipe validation, build and
(later) tests. & TD-03, TD-07 & Cheapest, highest-leverage safety net. \\
RN-02 & Add golden-file tests for the core rendering engine. & TD-01, Chapter~\ref{ch:testing} &
Protects the shared engine and all templates at once. \\
RN-03 & Ship a NativeWind/Tailwind theme defining the shadcn tokens
(\texttt{background}, \texttt{foreground}, \texttt{border}, \texttt{ring},
\dots). & KI-01 & Without it, generated shadcn screens may not style correctly. \\
RN-04 & Declare \texttt{expo-blur} on the three recipes that offer a glass
variant. & KI-02 & One-line fixes that make the glass variants functional. \\
RN-05 & Make generated auth hooks obviously stubbed and non-logging: remove
credential \texttt{console.log}, and mark the hook as a placeholder in its
comment. & KI-03, KI-11 & Reduces the chance of shipping non-functional auth. \\
RN-06 & Add a root \texttt{README} and a contribution guide. & TD-09, Chapter~\ref{ch:contribution-guide} &
Onboarding and discoverability. \\
RN-07 & Canonicalise the repository/brand identity across all metadata. &
SEO-04, GEO-01/02/05 & Removes citation ambiguity. \\
RN-08 & Add extension packaging to CI (package on tag; optionally publish or
attach the VSIX). & RM-08 & Removes the manual release step and the reliance on
local builds. \\
\bottomrule
\end{longtable}

\section{Mid-term recommendations}

\begin{longtable}{p{1.4cm}p{5.4cm}p{8.4cm}}
\caption{Mid-term recommendations}\\
\toprule
\textbf{ID} & \textbf{Recommendation} & \textbf{Rationale} \\
\midrule

RM-10 & Extract the shared orchestration layer (detection, option resolution,
primitives, file writing, manifest) into the core package. & Removes the highest
structural debt (TD-08); both clients call one implementation. \\
RM-11 & Replace hand-written validation with a real JSON Schema validator, or
generate validation from the schema. & Makes \texttt{schema.json} the single
source of truth (TD-07). \\
RM-12 & Implement one of the advertised-but-missing options end to end (Tamagui
is the largest and most requested). & Closes a configuration inconsistency
(KI-04, RM-03). \\
RM-13 & Introduce a compiled-template cache. & Small, safe performance win
(TD-05, PERF-01). \\
RM-14 & Harden the registry trust boundary: path containment, dependency-name
validation, HTTPS enforcement and content hashes. & Addresses high-severity
security findings (SEC-05, SEC-08, SEC-09, SEC-10, SEC-11). \\
RM-15 & Add dependency-update automation and an audit step. & Supply-chain hygiene
(TD-14). \\
\bottomrule
\end{longtable}

\section{Long-term / architectural directions}

\textbf{Status:} RECOMMENDED (strategic; not required for the current product).

\begin{itemize}
  \item \textbf{Remote-first catalog option.} If the catalog grows well beyond
        today's size, offer a client that fetches an index and loads recipes on
        demand, keeping the full inline bundle as an offline fallback
        (Chapter~\ref{ch:scalability}, SCALE-01/02).
  \item \textbf{A human-readable catalog site.} Derive a browsable page from the
        registry index to serve discoverability (SEO-05) and answer engines
        (Chapter~\ref{ch:aeo}).
  \item \textbf{Recipe authoring tooling.} A generator and preview for new
        recipes would lower the barrier to expanding the catalog (RM-06).
  \item \textbf{Cross-language reach.} The engine is TypeScript/Handlebars bound;
        a documented registry contract already exists, so other front-ends could
        consume it.
\end{itemize}

\section{Sequencing}

\begin{table}[H]
\centering
\caption{Suggested sequencing}
\begin{tabular}{p{5.6cm}p{4.2cm}p{4.4cm}}
\toprule
\textbf{Horizon} & \textbf{Focus} & \textbf{Items} \\
\midrule
Immediate & Correctness and safety nets & RN-01, RN-02, RN-03, RN-04, RN-05 \\
Short & Onboarding and consistency & RN-06, RN-07, RN-08, TD-10 \\
Mid & Architecture and features & RM-10, RM-11, RM-12, RM-13 \\
Long & Reach and hardening & RM-14, RM-15, remote-first catalog, catalog site \\
\bottomrule
\end{tabular}
\end{table}

\section{Summary}

The repository's own roadmap is small and mostly honest, marred only by listing
two already-implemented CLI features as outstanding. The highest-value work is
not new features but the verification and correctness items above: a CI gate, a
golden-file test suite, and the three output-correctness fixes that protect the
product's core promise.
